% GENERATED FILE. Edit canonical lessons, then run npm run generate:books.
% canonical-source-hash: fnv1a64:e40e7f6a3b4ad73d
% canonical-lessons: SA-C346-aksepah, SA-C346-khandayati, SA-C346-nirasyati, SA-C346-samsayah, SA-C346-vastutah

\chapter{Saying No to a Claim}
\label{ch:sa-saying-no-to-a-claim}
\hlchaptermodality{346}

\begin{chapteropening}
\textbf{By the end of this chapter:} \emph{I can object to a claim and reject it.}

Complete the last lesson of chapter 346: I can object to a claim and reject it.
\end{chapteropening}

\section[\texorpdfstring{\sk{आक्षेपः} (ākṣepaḥ)}{आक्षेपः (ākṣepaḥ)}]{\sk{आक्षेपः} (ākṣepaḥ) — an objection}

\label{lesson:SA-C346-aksepah}



\begin{quote}
Before the new one: say the Sanskrit for proves, then the Sanskrit for certainty.
\end{quote}

\subsection*{You'll want to know: \sk{आक्षेपः}}
\textbf{\sk{आक्षेपः}} — \emph{ākṣepaḥ} — \textquotedblleft{}an objection\textquotedblright{}.

\begin{quote}
\sk{कस्य} \sk{आक्षेपः} \sk{अस्ति}?
\end{quote}

\begin{quote}
\emph{kasya ākṣepaḥ asti?}
\end{quote}

\begin{quote}
Who has an objection?
\end{quote}

A masculine noun.

\begin{cousinweb}
\textbf{\sk{आक्षेपः}} (\emph{ākṣepaḥ}) is built on \emph{ā} and the root √kṣip, \textquotedblleft{}to throw\textquotedblright{}; MW gives \emph{ā-kṣepa} \textquotedblleft{}drawing together\textquotedblright{} and, among its later senses, \textquotedblleft{}objection\textquotedblright{} (MW).
\end{cousinweb}

\begin{grammarlens}[title={What you've built}]
One more word for giving reasons and opinions.
\end{grammarlens}

\subsection*{Your turn}
\begin{itemize}
\raggedright
  \item \emph{Say it:} \emph{ākṣepaḥ}
  \item \emph{Say it:} \emph{ākṣepaḥ}, once more
  \item \emph{Say it:} \emph{kasya ākṣepaḥ asti?} — Who has an objection?
  \item \emph{Recall:} say the Sanskrit for proves, then the Sanskrit for certainty, then say \emph{ākṣepaḥ} again
\end{itemize}

\begin{tcolorbox}[breakable,colback=teal!4,colframe=teal!35!black,title={Before you move on}]
What does \sk{आक्षेपः} mean? (\textquotedblleft{}An objection\textquotedblright{}.) Say it once more.
\end{tcolorbox}
\section[\texorpdfstring{\sk{खण्डयति} (khaṇḍayati)}{खण्डयति (khaṇḍayati)}]{\sk{खण्डयति} (khaṇḍayati) — refutes}

\label{lesson:SA-C346-khandayati}



\begin{quote}
Before the new one: say the Sanskrit for certainty, then the Sanskrit for an objection.
\end{quote}

\subsection*{You'll want to know: \sk{खण्डयति}}
\textbf{\sk{खण्डयति}} — \emph{khaṇḍayati} — \textquotedblleft{}refutes\textquotedblright{}.

\begin{quote}
\sk{पण्डितः} \sk{अरुणस्य} \sk{युक्तिं} \sk{खण्डयति।}
\end{quote}

\begin{quote}
\emph{paṇḍitaḥ aruṇasya yuktiṃ khaṇḍayati.}
\end{quote}

\begin{quote}
The scholar refutes Arun's argument.
\end{quote}

Past (he …): \sk{सः} \sk{युक्तिं} \sk{खण्डितवान्}, \emph{saḥ yuktiṃ khaṇḍitavān}.

\begin{cousinweb}
\textbf{\sk{खण्डयति}} (\emph{khaṇḍayati}) is built on \emph{khaṇḍa}, \textquotedblleft{}a piece, a part\textquotedblright{}; MW gives the verb the senses \textquotedblleft{}to break, to tear\textquotedblright{} and \textquotedblleft{}to refute\textquotedblright{} (MW).
\end{cousinweb}

\begin{grammarlens}[title={What you've built}]
One more word for giving reasons and opinions.
\end{grammarlens}

\subsection*{Your turn}
\begin{itemize}
\raggedright
  \item \emph{Say it:} \emph{khaṇḍayati}
  \item \emph{Say it:} \emph{khaṇḍayati}, once more
  \item \emph{Say it:} \emph{paṇḍitaḥ aruṇasya yuktiṃ khaṇḍayati.} — The scholar refutes Arun's argument.
  \item \emph{Recall:} say the Sanskrit for certainty, then the Sanskrit for an objection, then say \emph{khaṇḍayati} again
\end{itemize}

\begin{tcolorbox}[breakable,colback=teal!4,colframe=teal!35!black,title={Before you move on}]
What does \sk{खण्डयति} mean? (\textquotedblleft{}Refutes\textquotedblright{}.) Say it once more.
\end{tcolorbox}
\section[\texorpdfstring{\sk{निरस्यति} (nirasyati)}{निरस्यति (nirasyati)}]{\sk{निरस्यति} (nirasyati) — rejects (a claim)}

\label{lesson:SA-C346-nirasyati}



\begin{quote}
Before the new one: say the Sanskrit for an objection, then the Sanskrit for refutes.
\end{quote}

\subsection*{You'll want to know: \sk{निरस्यति}}
\textbf{\sk{निरस्यति}} — \emph{nirasyati} — \textquotedblleft{}rejects (a claim)\textquotedblright{}.

\begin{quote}
\sk{अध्यापकः} \sk{मम} \sk{उत्तरं} \sk{निरस्यति।}
\end{quote}

\begin{quote}
\emph{adhyāpakaḥ mama uttaraṃ nirasyati.}
\end{quote}

\begin{quote}
The teacher rejects my answer.
\end{quote}

Past (he …): \sk{सः} \sk{मम} \sk{उत्तरं} \sk{निरस्तवान्}, \emph{saḥ mama uttaraṃ nirastavān}.

\begin{cousinweb}
\textbf{\sk{निरस्यति}} (\emph{nirasyati}) is built on \emph{nir}, \textquotedblleft{}out\textquotedblright{}, and the root √as, \textquotedblleft{}to throw\textquotedblright{}; MW gives \emph{nir-√as} \textquotedblleft{}to cast out\textquotedblright{} and \textquotedblleft{}to reject, to refuse\textquotedblright{}, and \emph{nir-asta} \textquotedblleft{}rejected, refuted\textquotedblright{} (MW).
\end{cousinweb}

\begin{grammarlens}[title={What you've built}]
One more word for giving reasons and opinions.
\end{grammarlens}

\subsection*{Your turn}
\begin{itemize}
\raggedright
  \item \emph{Say it:} \emph{nirasyati}
  \item \emph{Say it:} \emph{nirasyati}, once more
  \item \emph{Say it:} \emph{adhyāpakaḥ mama uttaraṃ nirasyati.} — The teacher rejects my answer.
  \item \emph{Recall:} say the Sanskrit for an objection, then the Sanskrit for refutes, then say \emph{nirasyati} again
\end{itemize}

\begin{tcolorbox}[breakable,colback=teal!4,colframe=teal!35!black,title={Before you move on}]
What does \sk{निरस्यति} mean? (\textquotedblleft{}Rejects (a claim)\textquotedblright{}.) Say it once more.
\end{tcolorbox}
\section[\texorpdfstring{\sk{संशयः} (saṃśayaḥ)}{संशयः (saṃśayaḥ)}]{\sk{संशयः} (saṃśayaḥ) — a doubt}

\label{lesson:SA-C346-samsayah}



\begin{quote}
Before the new one: say the Sanskrit for refutes, then the Sanskrit for rejects (a claim).
\end{quote}

\subsection*{You'll want to know: \sk{संशयः}}
\textbf{\sk{संशयः}} — \emph{saṃśayaḥ} — \textquotedblleft{}a doubt\textquotedblright{}.

\begin{quote}
\sk{अत्र} \sk{संशयः} \sk{न} \sk{अस्ति।}
\end{quote}

\begin{quote}
\emph{atra saṃśayaḥ na asti.}
\end{quote}

\begin{quote}
There is no doubt about it.
\end{quote}

A masculine noun.

\begin{cousinweb}
\textbf{\sk{संशयः}} (\emph{saṃśayaḥ}) is built on \emph{sam} and the root √śī, \textquotedblleft{}to lie, to rest\textquotedblright{}; MW gives \emph{saṃ-śaya} \textquotedblleft{}uncertainty, irresolution, hesitation, doubt\textquotedblright{} (MW).
\end{cousinweb}

\begin{grammarlens}[title={What you've built}]
One more word for giving reasons and opinions.
\end{grammarlens}

\subsection*{Your turn}
\begin{itemize}
\raggedright
  \item \emph{Say it:} \emph{saṃśayaḥ}
  \item \emph{Say it:} \emph{saṃśayaḥ}, once more
  \item \emph{Say it:} \emph{atra saṃśayaḥ na asti.} — There is no doubt about it.
  \item \emph{Recall:} say the Sanskrit for refutes, then the Sanskrit for rejects (a claim), then say \emph{saṃśayaḥ} again
\end{itemize}

\begin{tcolorbox}[breakable,colback=teal!4,colframe=teal!35!black,title={Before you move on}]
What does \sk{संशयः} mean? (\textquotedblleft{}A doubt\textquotedblright{}.) Say it once more.
\end{tcolorbox}
\section[\texorpdfstring{\sk{वस्तुतः} (vastutaḥ)}{वस्तुतः (vastutaḥ)}]{\sk{वस्तुतः} (vastutaḥ) — in fact; really}

\label{lesson:SA-C346-vastutah}



\begin{quote}
Before the new one: say the Sanskrit for rejects (a claim), then the Sanskrit for a doubt.
\end{quote}

\subsection*{You'll want to know: \sk{वस्तुतः}}
\textbf{\sk{वस्तुतः}} — \emph{vastutaḥ} — \textquotedblleft{}in fact; really\textquotedblright{}.

\begin{quote}
\sk{वस्तुतः} \sk{सः} \sk{दरिद्रः} \sk{अस्ति।}
\end{quote}

\begin{quote}
\emph{vastutaḥ saḥ daridraḥ asti.}
\end{quote}

\begin{quote}
In fact he is poor.
\end{quote}

\begin{cousinweb}
\textbf{\sk{वस्तुतः}} (\emph{vastutaḥ}) is built on \emph{vastu}, \textquotedblleft{}a thing, the real\textquotedblright{}, with the ending \emph{-tas}; MW gives \emph{vastu-tas} \textquotedblleft{}in fact, in reality, actually\textquotedblright{} (MW).
\end{cousinweb}

\begin{grammarlens}[title={What you've built}]
One more word for giving reasons and opinions.
\end{grammarlens}

\subsection*{Your turn}
\begin{itemize}
\raggedright
  \item \emph{Say it:} \emph{vastutaḥ}
  \item \emph{Say it:} \emph{vastutaḥ}, once more
  \item \emph{Say it:} \emph{vastutaḥ saḥ daridraḥ asti.} — In fact he is poor.
  \item \emph{Recall:} say the Sanskrit for rejects (a claim), then the Sanskrit for a doubt, then say \emph{vastutaḥ} again
\end{itemize}

\begin{tcolorbox}[breakable,colback=teal!4,colframe=teal!35!black,title={Before you move on}]
What does \sk{वस्तुतः} mean? (\textquotedblleft{}In fact; really\textquotedblright{}.) Say it once more.
\end{tcolorbox}
